\section{Experimental protocol}
\label{sec:protocol}
% Source of everything in this section: docs/plan/two-track-plan.md sections 2.2-2.6 (the
% current plan), and docs/paper/score-spec.md section 2 (samplers). None of these
% experiments has been run at full scale; sizes and thresholds are the planned ones and may
% be re-sized at gate G1 (pilot audit), as the plan allows.

\subsection{Audit subjects: the E0 family}
\label{sec:protocol-e0}

The scores are audited on checkpoints of domain-randomised (DR) PPO agents trained on the
full generator (\cref{tab:experiments}). E0 is the main subject. Three manipulations
separate the factors RQ2 asks about: E0-sym replaces E0's privileged critic with the
symmetric one standard in UED and is co-primary; E0-const fixes the learning rate and the
entropy coefficient, so that a change with training phase cannot be a change of schedule;
and E0-short halves the preview distance $d_{\max}$, which changes what the agent can
observe.
% src: docs/plan/two-track-plan.md section 2.4 (experiments table)
Checkpoints are log-spaced from 10\,M to 1\,B steps. Training phases are defined by window
completion on A\textsubscript{DR}: below 50\%, 50--90\%, and 90\% or more.
% src: docs/plan/two-track-plan.md section 2.3 ("All are reported by phase ...")

\begin{table}[t]
\centering
\small
\caption{Planned experiments. Seed counts for C are re-sized at the pilot gate G1 from
E0's variance. A\textsubscript{QSS} is secondary and is not used for any claim in this
paper.}
\label{tab:experiments}
% src: docs/plan/two-track-plan.md section 2.4 (table "Experiments")
\begin{tabularx}{\linewidth}{@{}l>{\raggedright\arraybackslash}X>{\raggedright\arraybackslash}p{3.7cm}l@{}}
\toprule
ID & What & Size & Answers \\
\midrule
E0 & DR, asymmetric critic & 5 seeds $\times$ 1\,B steps; 16 checkpoints & RQ1, RQ2 \\
E0-sym & DR, symmetric critic (co-primary) & 3 $\times$ 1\,B; 16 checkpoints & critic privilege \\
E0-const & constant learning rate and entropy coefficient & 2 $\times$ 1\,B; 8 checkpoints & schedule vs.\ phase \\
E0-short & preview distance $d_{\max}$ halved & 3 $\times$ 1\,B; 8 checkpoints & observability \\
E1 & fresh-score audit of every E0-family checkpoint & 168 checkpoints $\times$ \num{1412}
  levels $\times$ 32--64 windows & RQ1, RQ2 \\
E2 & identities \eqref{eq:pvl-identity}, \eqref{eq:maxmc-identity}; tie fraction & logged during E1 & RQ2 \\
C & curricula on the \num{4000}-level pool (\cref{sec:protocol-c}) & 21 runs $\times$ 1\,B;
  buffer snapshots at 8 checkpoints & in-buffer RQ1; RQ3 \\
E5 & JaxUED maze, JaxUED's and SFL's own code at defaults & 5 seeds $\times$ 16 checkpoints
  $\times$ \num{2000} mazes & external check \\
A\textsubscript{QSS} & bias of quasi-steady-state labels & \num{20000} levels, 1 seed, 8
  checkpoints, 4 windows & secondary \\
\bottomrule
\end{tabularx}
\end{table}
% 168 = 5 x 16 + 3 x 16 + 2 x 8 + 3 x 8 (E0, E0-sym, E0-const, E0-short checkpoints).
\VERIFY{the plan gives E1 as \num{1412} levels per checkpoint, but A\textsubscript{DR} +
A\textsubscript{tail} + A\textsubscript{OOD} = \num{1400}; reconcile the extra 12 (they
must not be real circuits)}
\TODO{the asymmetric critic for E0 is not implemented yet (\cref{sec:env-mdp})}

\subsection{Targets and windows}
\label{sec:protocol-targets}

The headline targets are absolute regret $\Reg = \Jstar - \Jpi$ and normalised regret
$\Reg/\Jstar$, both primary. $\Jpi$ is the expected return of the \emph{stochastic}
policy, as UED's samplers see it; the deterministic policy's return is secondary.
% src: docs/plan/two-track-plan.md section 2.2 ("Episode", "Targets")
For each level and checkpoint we roll out 32 stochastic windows, raised adaptively to 64
when $0.05 < \hat p < 0.95$. Windows 1--4 feed the operational scores and windows 5 onward
the target, so a score and its target never share a rollout; high-budget scores are
cross-fitted on odd and even windows.
% src: docs/plan/two-track-plan.md section 2.2 ("Window assignment")
Two secondary targets remove known confounds: distance lost in 60\,s against the
undiscounted maximum-distance optimum on 600 A\textsubscript{DR} levels, which removes the
discount weighting; and a hurdle split into crash probability and regret given a
completed window. Total regret stays the headline because it is what UED optimises.
% src: docs/plan/two-track-plan.md section 2.2 ("Secondary targets")

\paragraph{Label noise} Two hundred pilot levels get at least three extra starts, which
gives a label standard deviation. A late-phase degradation is claimed only where the
median regret is at least three times that standard deviation, and we report
reliability-corrected correlations alongside raw ones.
% src: docs/plan/two-track-plan.md section 2.2 ("Label noise")

\subsection{Curricula on one labelled pool (C)}
\label{sec:protocol-c}

Every curriculum arm trains on the same \num{4000}-level pool, disjoint from the audit sets,
so all arms can be scored against the same labels and the in-buffer audit sees real
buffers. The arms are DR on the pool (5 seeds), \PLRp{} with PVL (5), \PLRp{} with
oracle-latest (5), \PLRp{} with MaxMC as implemented (3), and SFL (3); E0 serves as DR on
the full generator. Each runs for 1\,B steps with buffer snapshots at 8 checkpoints, and is
evaluated on A\textsubscript{DR}, A\textsubscript{tail} and A\textsubscript{OOD}.
% src: docs/plan/two-track-plan.md section 2.4 (row C; "Level sets")
The oracle arm scores a write group by $\Jstar - \frac{1}{E}\sum_e G^\gamma_e$, the
oracle-latest definition of \citet{abdelsadek2025goal} (in their code, \code{oracle-actor});
unlike theirs, our score never includes a trailing partial episode.
% src: docs/paper/score-spec.md section 2.2 ("An oracle-latest arm ..."); docs/paper/sources-recent.md (c)

\paragraph{Racing \PLRp{}} We keep JaxUED's \code{LevelSampler}: rank prioritisation with
temperature $\beta$, a staleness mix $\rho$, admission into the slot of lowest combined
weight only if the new score is strictly higher, a duplicate check, and gradient steps
only on replayed levels (\cref{tab:plr}).
% src: docs/paper/score-spec.md section 2.1; docs/paper/sources-jaxued.md section 2
Our windows outlast a training chunk, so the sampler works per slot and per chunk:
replay assignments are pre-drawn from the buffer state at the chunk start, a finished
window writes its score at the end of the chunk, a new level gets one full exploratory
window before it can enter the buffer, and the PPO loss is masked to replayed transitions
(\dev{15}--\dev{19}). Scores are accumulated across chunks with GAE bootstrapped at each
chunk end and written when the window ends; the audit's operational score calls the same
function. Unlike JaxUED, whose score is $-\infty$ when no episode completes within the
rollout, no level can be scored $-\infty$ here.
% src: docs/paper/score-spec.md section 2.2; docs/plan/two-track-plan.md section 2.2
%      ("Racing-PLR scoring protocol")
With $\beta = 0.3$ and \num{1000} slots, the top-ranked level carries 0.872 of the score
weight and, after the staleness mix, at least 0.61 of all replays. This is the reference
behaviour and C replicates it.
% src: docs/paper/score-spec.md section 2.1 (paragraph after the parameter table; test P-9)

\begin{table}[t]
\centering
\small
\caption{\PLRp{} parameters: JaxUED's class defaults, its maze example's defaults, the
Robust PLR paper (MiniGrid), and ours.}
\label{tab:plr}
% src: docs/paper/score-spec.md section 2.1 (parameter table)
\begin{tabular}{@{}lllll@{}}
\toprule
Parameter & Class default & Maze example & Robust PLR & Ours (C) \\
\midrule
Capacity & required & 4000 & 4000 & 1000 (\dev{20}) \\
Replay probability & 0.95 & 0.8 & 0.5 & 0.8 \\
Staleness $\rho$ & 0.5 & 0.3 & 0.3 & 0.3 \\
Prioritisation / $\beta$ & rank / 1.0 & rank / 0.3 & rank / 0.3 & rank / 0.3 \\
Minimum fill ratio & 1.0 & 0.5 & -- & 0.5 \\
Duplicate check & no & yes & -- & yes, by level id \\
Exploratory gradient updates & -- & no (\PLRp{}) & -- & no \\
$\gamma$ / $\lambda$ & -- & 0.995 / 0.98 & 0.995 / 0.95 & 0.995 / 0.95 \\
\bottomrule
\end{tabular}
\end{table}
\VERIFY{our minimum fill ratio of 0.5 is taken from the maze example's defaults and
awaits sign-off (spec section 2.1)}

\paragraph{Racing SFL} Every 50 chunks the parameters are frozen, the whole pool is
shuffled, each level is scored by $\hat p(1-\hat p)$ from four independent windows (S1),
and the top \num{1000} form the buffer; between refreshes, half of the environment slots
draw uniformly from the buffer and half from the pool, and every transition is trained on.
The shuffle restores the reference's random tie-breaking, and a NaN score raises an error
where the reference would silently select it. A refresh costs 19.2\,M scoring steps
($4000 \times 4 \times 1200$). The spec sets this against 13.1\,M training steps between
refreshes ($50 \times 2048 \times 128$, a ratio of 1.5$\times$; the reference ratios are
7.6$\times$ on MiniGrid and 0.76$\times$ on JaxNav), but that figure assumes 2,048
environments and 128-step chunks. With the plan's E0 configuration of 8,192 environments
$\times$ 64 steps it would be 26.2\,M training steps (0.73$\times$).
% src: docs/paper/score-spec.md sections 2.3, 2.4; 8,192 x 64: docs/plan/two-track-plan.md
%      section 4 (K4 row) and section 9 ("GPU configuration"); 50 x 8192 x 64 = 26,214,400;
%      19.2M / 26.2M = 0.73
\VERIFY{the SFL parameters $N$, $M$ and $T$ are proposals in the spec (section 6) and are
not fixed in the plan; the training-step count per refresh depends on the open chunk
length and environment count (\cref{sec:est-notation})}

\subsection{External check: JaxUED's maze (E5)}
\label{sec:protocol-e5}

To check that the audit's conclusions are not specific to racing, E5 repeats it in
JaxUED's 13$\times$13 maze with JaxUED's own DR and PLR-MaxMC code and the SFL authors'
MiniGrid code, all at their defaults, with agent views of 5$\times$5 and 13$\times$13
and a buffer audit as well as the fresh audit.
% src: docs/plan/two-track-plan.md section 2.4 (row E5); docs/paper/sources-sfl.md section 1.3
\VERIFY{the plan says ``JaxUED's own DR, PLR-MaxMC and SFL code''; the docs show no SFL
implementation in JaxUED, so we read SFL as the SFL authors' \code{minigrid\_sfl.py}, which
wraps JaxUED's maze}
The optimum is exact: a breadth-first search over (position, direction) finds the fewest
actions $k$ that end by stepping into the goal, so $J^* = \gamma^{k-1}(1 - 0.9\,k/250)$, or 0
if the goal is unreachable within 250 steps.
% src: src/apexzero/maze/oracle.py (module docstring)
We do not use JaxUED's own \code{optimal\_value}, whose time penalty and discount appear to
be one step off under the environment's reward timing.
% src: docs/paper/sources-jaxued.md section 3.1 ("Possible off-by-one ... my derivation; I did not run it")
\VERIFY{the off-by-one in JaxUED's \code{optimal\_value} is a derivation that has not been
confirmed by running it; confirm or drop this sentence}
Executing the search's plan in JaxUED's environment achieves exactly the computed optimal
return on every solvable random maze we tested, and JaxUED's maze example runs unmodified
from a vendored copy.
% src: tests/test_maze_oracle.py::test_bfs_optimum_is_achieved_exactly_in_the_jaxued_env
%      (atol 1e-6; at least 15 of 30 random mazes are solvable and checked);
%      docs/STATUS.md (P harness row); third_party/jaxued_examples/

\subsection{Endpoints and statistics}
\label{sec:protocol-endpoints}

\paragraph{Co-primary endpoints} With Holm correction within the family (to be listed in
\code{protocol-v1}; see below), each reported by phase and by critic type:
\begin{itemize}
  \item Spearman $\rho$ between the score and absolute regret, and between the score and
        normalised regret;
  \item RC@10\%, the regret captured by the score's top 10\% of levels relative to the
        most any 10\% could capture,
        $\mathrm{RC@10\%} = \sum_{\theta \in \mathrm{top}_{10}(S)} \Reg^{+}(\theta) \big/
        \sum_{\theta \in \mathrm{top}_{10}(\Reg)} \Reg^{+}(\theta)$, with
        $\Reg^{+} = \max(\widehat{\Reg}, 0)$, so a perfect ranking scores 1 and a random
        one scores $0.1\sum_\theta \Reg^{+} / \sum_{\mathrm{top}_{10}(\Reg)} \Reg^{+}$ in
        expectation, which we report alongside;
  \item in-buffer $\rho$, between the stored score and current regret, with score age as
        a covariate;
  \item admission AUROC, how well the score separates levels whose regret is in the top
        quantile;
  \item for RQ3, the interquartile mean of held-out regret at 1\,B steps.
\end{itemize}
% src: docs/plan/two-track-plan.md section 2.3 ("Co-primary endpoints")
\paragraph{Negative regret estimates} A per-level estimate $\widehat{\Reg} = \Jstar -
\Jhat$ can be negative, through Monte Carlo noise in $\Jhat$ or a label the agent beats.
Rank correlations use the raw estimates; RC@10\% clips them at zero as above; and levels
on which an observed agent return beats $\Jstar$ stay in the endpoints, flagged as label
shortfalls and counted in the report of \cref{sec:answer-lb}.
\VERIFY{the plan defines RC@10\% only in words (``the share of regret captured by the
score's top 10\%''); the normalised form above and the handling of negative estimates are
our proposal and must be signed off in \code{protocol-v1}, together with the admission
quantile, which the plan leaves unspecified. Under this definition the provisional
thresholds below should be checked against the random-score expectation, which depends on
how concentrated regret is}
\TODO{list the Holm family in \code{protocol-v1}: which of the seven estimator rows,
the five co-primary endpoints, the three phases and the two critics (E0, E0-sym) form one
family, and whether absolute and normalised regret, and S1 and S2 learnability, count as
separate tests}

\paragraph{Diagnostics} A two-loading regression
$\mathrm{rank}(S) \sim a\,\mathrm{rank}(\Jstar) + b\,\mathrm{rank}(\Jhat)$, for which a
faithful score has $a > 0$ and $b < 0$; partial $\rho$; precision and contamination at
quantile $q$; a statistic $\eta$ computed from a simulated \PLRp{} replay process
(buffer, admission rule, staleness mix, $\beta \in \{0.1, 0.3, 1\}$); and lagged
fidelity, the score at checkpoint $k-1$ against regret at checkpoint $k$. We do not use
a static rank$^{-1/0.3}$ weighting, which puts 87\% of its mass on one level.
% src: docs/plan/two-track-plan.md section 2.3 ("Diagnostics")
\VERIFY{the plan names $\eta$ without defining it; we intend the share of simulated replays
that go to levels in the true top regret quantile, to be fixed in \code{audit-v1}}

\paragraph{Thresholds} Provisionally, a score is \emph{faithful} if $\rho \ge 0.5$ and
RC@10\% $\ge 0.6$, and \emph{uninformative} if $\rho \le 0.2$, or RC@10\% $\le 0.2$, or it
does not beat the best agent-only and policy-agnostic baseline (paired interval). The
final thresholds are set at G1 from the pilot's minimum detectable effect.
% src: docs/plan/two-track-plan.md section 2.3 ("Thresholds")

\paragraph{Statistics} Intervals come from a two-way bootstrap over seeds and levels, or a
crossed-effects model on Fisher-$z$ correlations. The primary trend test is a $t$-test on
per-seed slopes. Equivalence is tested by TOST with a margin from the pre-computed minimum
detectable effect. The PPO return scale $\sigma$, the entropy coefficient and the
observation-normalisation statistics enter as covariates. Aggregates follow
rliable~\citep{agarwal2021rliable}.
% src: docs/plan/two-track-plan.md section 2.3 ("Statistics")

\subsection{Pre-registration}
\label{sec:prereg}

The scorer, the window assignment, the level partition and the fallback for an
early-versus-mid comparison are frozen in the tag \code{audit-v1}; the endpoints,
thresholds, predictions and the outcome table below are frozen in \code{protocol-v1}. Both
are tagged at G1, after the pilot audit and before any full-scale result. A claim ledger
written by the figure scripts ties every number in the paper to the script and data that
produced it.
% src: docs/plan/two-track-plan.md sections 2.2, 2.6 ("Forking paths"), 4 (lockbox rules), 5.2 (G1)
\TODO{neither tag exists yet; this section describes what they will contain}

\paragraph{Predictions}
\begin{enumerate}[(P1)]
  \item The ratio PVL/L1 tends to 0.5 as critic error shrinks, in mode B (in mode F the
        window-end artefact of \cref{sec:est-notation} keeps it below 0.5).
  \item MaxMC as implemented has a high $\rho$ against absolute regret and a $\rho$ near
        zero against normalised regret.
  \item The error of unit-consistent MaxMC is dominated by the ceiling gap on hard levels.
  \item In-buffer fidelity is lower than fresh fidelity.
\end{enumerate}
% src: docs/plan/two-track-plan.md section 2.5 ("Pre-registered predictions");
%      mode-B restriction from docs/paper/score-spec.md sections 0.3 and 6

\paragraph{Outcomes} \Cref{tab:outcomes} maps each possible outcome to the claim it
licenses. Rows B, D and E trigger a follow-up experiment (E4) only if they hold for PVL,
MaxMC or learnability at the operational budget after Holm correction, the label-noise
floor and the reliability correction; every E4 design includes both DR arms and uses the
audit's units.
% src: docs/plan/two-track-plan.md section 2.5 ("E4 trigger")

\begin{table}[t]
\centering
\small
\caption{Pre-registered outcome table. Each row's claim is licensed only if its test
passes.}
\label{tab:outcomes}
% src: docs/plan/two-track-plan.md section 2.5 (outcome table)
\begin{tabularx}{\linewidth}{@{}>{\raggedright\arraybackslash}p{2.0cm}>{\raggedright\arraybackslash}X>{\raggedright\arraybackslash}X>{\raggedright\arraybackslash}p{2.3cm}@{}}
\toprule
Row & Test & Claim & Follow-up \\
\midrule
A, faithful & faithful in every phase, for both critics, fresh and in-buffer; TOST passes &
  ``With dense reward, scores track regret in continuous control'', scoped by E0-sym,
  E0-short and E5 & none \\
B, degrades late & $\rho_{\text{late}} - \rho_{\text{early}} \le -0.1$ with the interval
  excluding 0, still present in E0-const & ``Scores guide the curriculum only while the
  agent is incompetent, because of [E2 term]'' & E4-switch, with a $\Jhund$-scored control
  arm \\
C, uninformative & uninformative throughout & ``Ground truth confirms No Regrets in
  continuous control'' & none \\
D, track type & $\Delta\rho \ge 0.2$ between track-type strata & ``Score error is
  structured, so samplers under-replay [type]'' & E4-type \\
E, faithful but not useful & oracle-latest $\approx$ DR-pool in C & ``Regret is the wrong
  target here'' & E4-hybrid \\
F, manipulations disagree & symmetric vs.\ asymmetric critic, full vs.\ short preview, or
  racing vs.\ maze differ by $\ge 0.2$ & ``Fidelity depends on [critic privilege /
  observability]'' & none \\
G, fresh $\ne$ buffer & in-buffer $\rho$ differs from fresh $\rho$ by $\ge 0.2$ &
  ``DR-style audits misstate what curricula actually see'' & none \\
\bottomrule
\end{tabularx}
\end{table}

\paragraph{Strata and lockbox} Track-type strata are tagged from the tokens in the first
$\approx$1.4\,km of a strip, which carries 90\% of the discounted weight. No audit or
curriculum experiment runs on a real circuit. The twelve development layouts supply the
generator's statistics and hosted the development checks (the specialist, the
minimum-time replay gate and the solver spike). No number in this paper is computed on
the twelve lockbox circuits (\cref{sec:threats}).
% src: docs/plan/two-track-plan.md sections 2.4 ("Level sets") and 4 ("Lockbox rules")

\paragraph{Compute} The plan's total GPU budget, which also covers tuning and a companion
project, is 62--127 GPU-hours on one RTX 4090 without E4; the labels are budgeted at
1,200--3,000 CPU-hours, unverified until the TO-OCP spike finishes.
% src: docs/plan/two-track-plan.md section 8 ("Compute")
\RESULT{actual GPU-hours and CPU-hours used}
